\title{Large Language Models Can Automatically Engineer Features for Few-Shot Tabular Learning}

\begin{abstract}
Large Language Models (LLMs) have exhibited extraordinary proficiency in addressing complex and previously unencountered reasoning tasks, underscoring their significant promise for tabular learning—a capability fundamental to numerous practical scenarios. This paper introduces an innovative in-context learning paradigm, \textsf{FeatLLM}, which harnesses LLMs as feature engineers to synthesize transformed datasets tailored for optimal tabular prediction performance. The features constructed by this approach are subsequently leveraged by a straightforward predictive model, such as linear regression, achieving highly effective few-shot learning outcomes. Notably, the \textsf{FeatLLM} approach confines inference to utilizing the identified features alongside the simple downstream model, dispensing with LLM involvement during test time for each prediction. In contrast to prevailing LLM-centric strategies, \textsf{FeatLLM} obviates the need for per-instance LLM queries during inference. Furthermore, it requires only API-level interaction with LLMs and circumvents constraints related to prompt size. Experimental results spanning a diverse array of tabular datasets illustrate that \textsf{FeatLLM} discovers high-fidelity feature rules, substantially surpassing competing solutions, including TabLLM and STUNT, by an average of 10\%.
\end{abstract}

\section{Introduction}
Large Language Models (LLMs) have, in recent times, demonstrated a remarkable ability to provide relevant outputs for previously unseen tasks without the necessity for fine-tuning that is specific to a task~\cite{creswell2022selection,imani2023mathprompter,taylor2022galactica}. Through pretraining on vast and heterogeneous textual datasets, LLMs accumulate extensive background knowledge that can supplant specialized human expertise in various disciplines~\cite{Genomes2021,singhal2023large,wilcox2003role}. This accumulated knowledge enables the automation of a broad spectrum of applications, such as dialogue interpretation~\cite{finch2023leveraging} and codebase repair~\cite{fan2023automated}. Particularly noteworthy is that LLMs are capable of assimilating a given task’s context when prompted with succinct task descriptions and a limited set of examples, thereby identifying salient features and demonstrative samples~\cite{brown2020language,zhao2023survey}. This feature selection capacity accentuates the potential of LLMs to function as automated, expert-level feature engineers.

Efforts to employ the latent knowledge and reasoning strengths of LLMs have progressed into tabular learning. For example, prior knowledge indicating age-related elevation of certain health risks is advantageous for predictive modeling in medicine. Existing studies frequently formalize tasks and describe features through natural language, serialize tabular data accordingly, and input this representation into an LLM~\cite{dinh2022lift,wang2023anypredict}. These workflows commonly utilize either in-context learning methods~\cite{nam2023semi} or opt for parameter-efficient model fine-tuning~\cite{dinh2022lift,hegselmann2023tabllm}. The integration of LLM-derived background knowledge has consistently demonstrated superior performance, particularly outstripping traditional tabular learning baselines within limited data regimes~\cite{hegselmann2023tabllm}.

Existing approaches leveraging LLMs for tabular data analysis face notable challenges. In end-to-end predictive tasks, each data point necessitates at least one separate LLM inference, resulting in high computational overhead. Achieving strong performance with these methods typically involves fine-tuning the LLM, which is impractical for many advanced LLMs lacking open-source training infrastructure or public tuning interfaces. This is particularly pressing as state-of-the-art models have increasingly restricted access to inference-only APIs~\cite{anil2023palm,openai2023gpt4}. Moreover, several current techniques struggle when prompts become lengthy~\cite{liu2023lost}; as the number of tabular features increases and the text representations expand, these approaches frequently become unusable or see diminished accuracy. Collectively, these issues limit the viability of LLM-centric tabular learning for practical deployment.

To address these concerns, we propose a shift in the role of LLMs from end-to-end predictors to tools for feature creation. Our method employs LLMs for feature engineering, harnessing their background knowledge more strategically. We present \textsf{FeatLLM}, a new in-context learning paradigm designed to elicit the underlying “criteria” that inform LLM decision-making. For example, in medical diagnosis tasks, the LLM can surface explicit decision rules that associate certain feature patterns with disease prediction. These criteria or rules then inform the construction of new, derived features, which can serve as substitutes for raw variables in subsequent modeling steps. This repositions the LLM as a generator of informative features, enabling simpler models downstream—thus improving inference times once features are finalized. The approach capitalizes on the LLM’s capability for in-context learning, facilitating feature generation simply by incorporating illustrative samples in the prompts, thereby removing the necessity for model training. Finally, borrowing strategies from ensemble techniques~\cite{breiman2001random}, such as feature bagging, enables effective management of prompt length and complexity.

\textsf{FeatLLM} organizes the input prompt for novel feature engineering into two main phases: (1) analyzing the given problem to discern how features interact with target variables, and (2) leveraging insights from step one, along with several illustrative examples, to construct explicit rules that effectively separate classes. This systematic reasoning framework steers LLMs away from considering irrelevant features during rule formulation. The generated rules are executed as code over the data, transforming each sample into a binary indicator that reflects compliance with each specific rule. Subsequently, a simple predictive model is trained to map these engineered features to class probabilities. To further increase robustness, \textsf{FeatLLM} adopts an ensemble scheme, repeatedly generating diverse features with the aid of feature bagging. By judiciously selecting the quantity of features and samples included during the bagging process, the method addresses the bottleneck posed by large prompt sizes. As \textsf{FeatLLM} does not require gradient-based model training, it incurs low inference overhead, making it suitable for deployment in practical, real-world settings.

Empirical evaluation is conducted on 13 tabular datasets under low-data conditions, where \textsf{FeatLLM} demonstrates consistently high and stable performance. Our results indicate that LLMs can effectively leverage both background knowledge and information present within the data, yielding salient and discriminative feature representations. Our approach surpasses state-of-the-art few-shot learning baselines over a range of scenarios. An anonymized GitHub repository hosting the source code is available at~\texttt{https://github.com/Sungwon-Han/FeatLLM}.

\section{Related Work}

\subsection{Few-Shot Learning with Tabular Data} The progress in few-shot learning has facilitated the extraction of broadly applicable feature representations from datasets, even when annotated samples are scarce~\cite{chen2018closer}. While the field was initially oriented toward image recognition problems, few-shot learning now exhibits substantial capacity for transfer across an array of data types, encompassing text, audio, and more recently, tabular data formats~\cite{majumder2022few,nam2023stunt,schick2021exploiting}. Nevertheless, deriving insights from limited labeled examples raises concerns regarding the inadvertent modeling of coincidental or non-generalizable patterns~\cite{bartlett2021deep}. This risk has prompted newer research directions that augment model learning using supplementary information sources or encode domain-specific priors that bias the inductive process in beneficial ways. Among common approaches, semi-supervised learning frameworks harness large collections of unannotated data—often applying sophisticated data augmentations to yield robust models~\cite{ucar2021subtab,yoon2020vime}. Alternatively, one can achieve beneficial initialization via unsupervised pre-training regimes, which do not depend on specific downstream task knowledge~\cite{somepalli2021saint}. Common unsupervised learning strategies incorporate tasks such as reconstructing masked or corrupted features~\cite{arik2021tabnet,majmundar2022met}, or applying augmentation techniques in the context of contrastive learning~\cite{bahri2022scarf,chen2020simple}. The intrinsic heterogeneity characteristic of tabular datasets, however, complicates the design of augmentation strategies with wide applicability, resulting in a general lack of efficacy for these methods in few-shot environments~\cite{nam2023stunt}.

Contemporary research has suggested pre-training on a diverse collection of tabular datasets unrelated to the end task, followed by transfer to target applications~\cite{zhang2023generative,levin2022transfer,zhu2023xtab}. As an example, TransTab~\cite{wang2022transtab} utilizes transformer architectures to infer semantic dependencies between table columns via their column names in addition to cell values. Through this mode of training, the resulting models possess transferable knowledge of a variety of tables and their constituent features, which can be exploited in few-shot or even zero-shot settings. Another line of work, exemplified by TabPFN~\cite{hollmann2022tabpfn}, generates synthetic datasets combining multiple statistical distributions to facilitate pre-training of Bayesian neural networks. After this broad-based pre-training, predictions for new tasks are obtained by jointly processing support and query examples without any task-specific fine-tuning. The inductive biases accrued through exposure to a wide variety of tables and distributions allow such models to outperform classical tree-based methods, demonstrating significant promise for data-scarce scenarios.

\subsection{Language-Interfaced Tabular Learning} A further innovative method for injecting prior knowledge is through leveraging large language models (LLMs)~\cite{anil2023palm,openai2023gpt4}. Having been trained on expansive and varied text datasets, LLMs possess notable cross-domain generalization abilities, routinely exhibiting competence on tasks not present in their training distribution~\cite{brown2020language}. Recent investigations have attempted to exploit the encoded prior knowledge within LLMs for tabular data problems. A typical practice is to linearize structured tables into textual form and feed these serialized inputs to LLMs through natural language prompts~\cite{dinh2022lift,hegselmann2023tabllm,wang2023anypredict}. By carefully structuring prompts to include not only tabular data but also detailed descriptions of features and tasks, it becomes possible for the model to infer complex task-feature interactions in a zero-shot or few-shot regime. State-of-the-art work explores either fine-tuning LLMs with resource-efficient methods such as LoRA~\cite{hu2021lora} or IA3~\cite{liu2022few}, or capitalizes on in-context learning, supplying a handful of curated demonstration examples to guide model predictions without the need for direct parameter updates~\cite{dinh2022lift,hegselmann2023tabllm,nam2023semi,slack2023tablet}.

While previous methods have demonstrated effectiveness in enhancing few-shot tabular learning, their reliance on always routing inference through the LLM limits practical deployment, especially in industrial settings. The objective of this work is to diverge from traditional black-box strategies where the LLM predicts outcomes for each input in a holistic, end-to-end manner. Rather, this approach leverages the LLM to deduce class-specific conditions or rule sets directly. \textsf{FeatLLM} subsequently transforms these LLM-derived rules into executable program code for feature construction. This procedure allows for inference that bypasses the LLM, instead using in-context learning to derive and encode rules from available few-shot data and pre-existing knowledge.

\section{Methods}
\textsf{FeatLLM} is constructed to identify ``rules"—namely, sets of attribute conditions that correspond to specific classes—using the small number of example instances provided, as illustrated in Figure~\ref{fig:main_model}. These rules, extracted by the LLM, are then systematically converted into binary features indicating which of them are satisfied by a sample. The resulting binary feature vector is fed into a weighted dot product using non-negative, trainable parameters to generate an estimate of class probabilities. During training, this allows the model to calibrate and assign significance to each generated feature based on its contribution to predicting class membership (see Section~\ref{Sec:training}). This methodology is iteratively executed across several subsamples using bagging, with results later merged through an ensemble strategy. The following sections articulate the formal problem statement and provide further specification for each component of the approach.

\vspace*{-0.12in}
\paragraph{Problem formulation.} Consider a tabular dataset containing $N$ labeled instances, denoted as $\mathcal{D} = \{(\mathbf{x}^i, \mathbf{y}^i)\}_{i=1}^{N}$. In this collection, each sample $\mathbf{x}^i$ is a $d$-dimensional vector representing the $d$ input attributes. Associated with these attributes are natural language feature names $F = \{f_j\}_{j=1}^{d}$, accompanied by their respective definitions provided externally. Within the framework of a classification problem, every label $\mathbf{y}^i$ belongs to a fixed set of classes. For $k$-shot learning scenarios, the model is trained on a subset of $k$ ($<N$) labeled samples chosen randomly from $\mathcal{D}$.

\subsection{Prompt Design for Extracting Rules}
\label{Sec:prompt_design}
To facilitate more precise and interpretable reasoning in rule extraction by LLMs, we steer the problem-solving process to align with the methodology an expert human would use in tabular data analysis. The prompts we use are structured into three principal sections (see Fig.~\ref{fig:prompt_for_rule} and Appendix~\ref{sec:example_rule}):

\vspace*{-0.12in}
\paragraph{Basic information description.} This segment outlines the foundational information necessary for addressing the problem (refer to the orange-highlighted portion in Fig.~\ref{fig:prompt_for_rule}). It contains a formal statement of the task along with its labeling scheme, detailed descriptions of the features, and illustrative examples. The description of the task is posed in the form of a question (e.g., ``Does this patient's myocardial infarction complications data indicate chronic heart failure? Yes or no?". Additional examples can be found in Appendix~\ref{sec:task_desc}). Each feature is described in terms of its type of value and, where applicable, its definition; categorical features may also be accompanied by a listing of valid categories. For examples, a handful of labeled training points are transformed into textual form and shown with their true labels. The transformation employs a serialization format adapted from prior literature~\cite{dinh2022lift,hegselmann2023tabllm}:

\begin{align}
&\text{Serialize}(\mathbf{x}^i,\mathbf{y}^i, F) = \nonumber \\ 
&``\ f_1\ \text{is}\ \mathbf{x}^i_1.\ ... \ f_d\  \text{is}\  \mathbf{x}^i_d.\ \ \text{Answer:}\  \mathbf{y}^i\ ”
\end{align}

\vspace*{-0.12in}
\paragraph{Reasoning instruction.} Our objective is to improve the LLM's reasoning ability by integrating explicit reasoning instructions into the prompt (highlighted as blue text in Fig.~\ref{fig:prompt_for_rule}). The reasoning instruction opens with a sentence akin to that used in the chain-of-thought methodology~\cite{wei2022chain}. Subsequently, we structure the inference process into two discrete stages. The first stage prompts the LLM to hypothesize causal associations or trends between the features and the specified task independently, without drawing upon provided example demonstrations; instead, it leverages general knowledge and commonsense reasoning, maximizing the use of its pre-existing knowledge about the domain. In the second stage, the LLM utilizes both the initial hypotheses and example demonstrations to formulate specific rules for each class label. This sequential approach guards against the selection of irrelevant patterns or spurious correlations, helping the model to prioritize features with meaningful impact. To balance between underfitting and excessive complexity, and to constrain prompt length, we stipulate a constant number of rules (specifically, 10)\footnote{Discussion of rule number choices is in Section~\ref{sec:analysis}.} to be generated for every class. Furthermore, to prevent type inconsistencies, we instruct the LLM to reference the type information for each feature provided in the basic information section when constructing rules.

\vspace*{-0.12in}
\paragraph{Response instruction.} To streamline rule extraction and downstream usage, we supply the LLM with targeted instructions on formatting its outputs (indicated as yellow text in Fig.~\ref{fig:prompt_for_rule}). The prompt clearly outlines both the structure required for each answer class (i.e., response formatting) and the specific syntax each rule should follow (i.e., rule formatting), allowing the LLM to select formats that align with feature types. If not further restricted, the LLM is permitted to synthesize multiple rules using logical operators such as AND or OR. A concrete example of the model's response appears in Appendix~\ref{sec:outcome-rule}.

\subsection{Inferring Class Likelihood via Rules}
\label{Sec:training}

\paragraph{Parsing rules for feature generation.} In this stage, we generate binary features from the rules produced earlier. Specifically, for each class, a binary indicator is assigned to reflect whether the sample meets the relevant rule criteria (represented as ‘0’ or ‘1’). These features contribute to predicting the class membership probability for each instance. Nevertheless, as the rules originate from LLM outputs in natural language, it is necessary to construct an effective parsing approach for systematic feature extraction. While we attempt to impose strict formatting on both answers and rules during rule synthesis to aid downstream parsing, outputs from the LLM may still display unpredictable or inconsistent syntax~\cite{kovavc2023large}. For example, the LLM may rephrase symbols such as “$>$” or “$<$” with their verbal counterparts, like “greater than” or “less than,” thereby increasing parsing complexity.

Rather than relying on intricate programmatic solutions to manage such noisy outputs, we utilize the LLM itself as the parsing agent. Owing to its proficiency in semantic interpretation regardless of syntactic form, and its coding capabilities learned from extensive code data, the LLM can be effectively prompted to convert textual rules into executable code. To harness these strengths, we construct prompts including details such as the target function’s name, descriptions of input and output, and the extracted rules, which are then supplied to the LLM. As illustrated in Figures~\ref{fig:prompt_for_parsing} and~\ref{sec:example_parsing}-\ref{sec:example_parse_outcome} in the Appendix, these prompts, along with the LLM’s generated responses, demonstrate the parsing procedure. The resulting code snippets are then run through Python’s \texttt{exec()} method to facilitate automated feature computation.

\vspace*{-0.12in}
\paragraph{Inferring class likelihood.} When new binary features are derived from the rules associated with each class, one straightforward approach for estimating the likelihood of a sample's class membership is to tally the number of rules that the sample meets for each class (i.e., computing the sum of its class-specific binary features). Yet, the relative significance of different rules and features can vary, so their weights should be learned from the training data. To address this, we utilize a standard machine learning (ML) model designed to determine the relative importance of each class's binary features. For instance, we can employ a bias-free linear model for this purpose. Define $\mathbf{z}_k^i \in \{0, 1\}^R$ as the vector of binary features produced for class $k$ from the input $\mathbf{x}^i$, where $R$ denotes the number of rules per class and $c$ is the number of total classes. The class probability vector $\mathbf{p}^i$ is computed via:

\begin{align}
\text{logit}_k^i &= \max(\mathbf{w}_k, 0) \cdot \mathbf{z}_k^i, \nonumber \\
\mathbf{p}^i &= \text{Softmax}({[\text{logit}_{1}^i, ..., \text{logit}_{c}^i]}). \label{eq:infer_prob}
\end{align}

Here, the vector $\mathbf{w}_k$ consists of trainable parameters in the linear model, quantifying each feature's importance. By restricting these weights to the non-negative domain, we guarantee that features extracted by the LLM cannot diminish a class's predicted probability during training. The logits are subsequently translated into normalized class probabilities via the softmax operation. To fit the parameters $\mathbf{w}_k$, cross-entropy loss is minimized using a small sample of labeled training data.

\vspace*{-0.12in}
\paragraph{Ensembling with bagging.} To improve predictive robustness, the procedure is conducted repeatedly to construct several inference models, combining their outputs through ensembling. With the set of model weights from $T$ independent runs, denoted as $\{\mathbf{w}[t]\}_{t=1}^T$, class probabilities for a new input are evaluated by averaging the outputs $\mathbf{p}^i$ (obtained as in Eq.~\ref{eq:infer_prob}) for each individual $\mathbf{w}[t]$.

To inject randomness into the rule-generation process for the ensemble, we employ a relatively high temperature setting during LLM inference. In addition, the order of few-shot examples is randomized, motivated by findings that the sequence of demonstrations can significantly affect LLM outputs~\cite{lu2022fantastically}\footnote{An exploration of rule diversity is presented in Appendix~\ref{sec:diversity}}. To exploit this diversity and improve predictive robustness, we implement bagging, selecting different subsets of either features or data instances for each ensemble iteration—a practice especially advantageous when dealing with large feature spaces or training sets. Our ensemble method confers two primary benefits. Firstly, if some of the rules generated by the LLM in specific trials are erroneous, correct rules identified in other rounds can offset them, ultimately strengthening the system’s resilience. This leverages the LLM's self-consistency property~\cite{wang2022self}, as rules that recur across independent trials are more likely to be reliable. Secondly, this ensemble technique helps overcome the constraint imposed by the finite prompt length in LLMs. For datasets whose tabular content exceeds the prompt limit, bagging effectively reduces input size, thereby maintaining robust model performance. Whereas a singular rule set might fail to represent the relationship among all features, ensembling a multitude of weak learners derived from several trials compensates for any single model’s limited coverage of features or instances.

\section{Experiments}
We assess \textsf{FeatLLM} on a range of tabular datasets, complemented by additional analyses to better understand the behavior and utility of the framework. Further results, including those concerning other LLM variants, qualitative observations, and supplementary experimental details, appear in the Appendix for brevity.

\subsection{Performance Evaluation}
\paragraph{Datasets.}
Our study leverages 13 datasets encompassing binary and multi-class classification scenarios: (1) Adult~\cite{asuncion2007uci}, tasked with determining if an individual’s annual income surpasses \$50,000; (2) Bank~\cite{moro2014data}, concerning prediction of a customer’s likelihood to purchase a term deposit; (3) Blood~\cite{yeh2009knowledge}, for identifying returning blood donors; (4) Car~\cite{kadra2021well}, which aims to assess car quality; (5) Communities~\cite{misc_communities_and_crime_183}, oriented toward forecasting regional crime rates; (6) Credit-g~\cite{kadra2021well}, distinguishing between high- and low-risk credit applicants; (7) Diabetes\footnote{\texttt{kaggle.com/datasets/uciml/pima-indians-diabetes-database}}, used for diagnosing diabetes in individuals; (8) Heart\footnote{\texttt{kaggle.com/datasets/fedesoriano/heart-failure-prediction}}, which forecasts coronary artery disease; and (9) Myocardial~\cite{misc_myocardial_infarction_complications_579}, for identifying cases of chronic heart failure.

In addition, we incorporate a set of newer datasets and artificially constructed datasets that have received less attention in the literature and were not present during the LLM pre-training phase: (10) Cultivars, used to forecast the grain yield potential for a particular soybean cultivar\footnote{\texttt{archive.ics.uci.edu/dataset/913}}; (11) NHANES, which tasks models with predicting an individual’s age group based on dataset records\footnote{\texttt{archive.ics.uci.edu/dataset/887}}; (12) Sequence-type, where the objective is to categorize sequences as belonging to one of four forms: arithmetic, geometric, Fibonacci, or Collatz; and (13) Solution-mix, which determines whether the concentration percentage from a blend of four solutions surpasses 0.5. The first two datasets were only recently made available on the UCI Machine Learning Repository post-September 2023, thereby assuring exclusion from the pre-training material used by the LLM API (gpt-3.5-turbo-0613) employed in our model. The latter two are generated synthetically, ensuring the LLM could not have had prior exposure and could not memorize their content, thus maintaining the integrity of the evaluation.

These datasets demonstrate wide variability in both their number of entries and overall difficulty, as outlined in Table~\ref{Tab:basic-desc}. Each dataset is accompanied by clear labeling and straightforward descriptions for individual attributes. For example, datasets such as `Communities' and `Myocardial' have more than 100 features, making them particularly challenging due to the limitations imposed by the LLM's available context window.

\vspace*{-0.12in}
\paragraph{Baselines.}
Our analysis includes nine benchmark methods for comparison. The initial three baselines represent established algorithms for tabular data modeling: (1) Logistic regression (LogReg), (2) XGBoost~\cite{chen2016xgboost}, and (3) RandomForest~\cite{ho1995random}. Two additional baselines assume accessibility to unlabeled samples, namely (4) SCARF~\cite{bahri2022scarf} and (5) STUNT~\cite{nam2023stunt}; however, acquiring abundant unlabeled data is often unrealistic in many real-world scenarios. (6) TabPFN~\cite{hollmann2022tabpfn}, a recent approach, produces artificial datasets with diverse distributions to enable model pre-training. The last group involves large language model-based baselines: (7) In-context learning (In-context)~\cite{wei2022emergent}, (8) TABLET~\cite{slack2023tablet}, and (9) TabLLM~\cite{hegselmann2023tabllm}. For in-context learning, several representative examples are included in the prompt without any modification to model parameters. TABLET augments prompts with extra cues, leveraging rules and prototypes from an auxiliary classifier to enrich inference. TabLLM makes use of parameter-efficient tuning, such as IA3~\cite{liu2022few}, applying few-shot updates to adapt the LLM for specific tabular learning tasks.

\vspace*{-0.12in}
\paragraph{Implementation details.}
\textsf{FeatLLM} relies on GPT-3.5 as its foundational LLM, but is intentionally constructed to remain flexible regarding the underlying LLM choice (experimental results utilizing alternative backbones can be found in Appendix~\ref{sec:backbones}). For LLM inference, the temperature parameter is adjusted to 0.5, and the top-p setting is kept at the default value of 1 as specified in the API. We specify the ensemble size to be 20 and fix the number of rules extracted at 10. The effects of different hyper-parameter selections are presented in Figure~\ref{fig:hyperparameter2} and discussed further in Appendix~\ref{sec:hyper-parameter}.
The linear model is trained using the Adam optimizer with a learning rate of 0.01 over 200 epochs. To select the optimal number of epochs, $k$-fold cross-validation is utilized. When benchmarking against established baselines, in-context learning is implemented using GPT-3.5, while TabLLM is evaluated with T0~\cite{sanh2022multitask} in line with its initial configuration. Since TabLLM limits the LLMs to those with open-source checkpoints, GPT-3.5 is excluded from comparison in this case. For classical machine learning models, such as LogReg, XGBoost, and RandomForest, hyper-parameters were identified through Grid-search in conjunction with $k$-fold cross-validation. The parameter $k$ is chosen to be either 2 or 4 to ensure the training set encompasses at least one instance from every class. Further specifics on the implementation can be found in Appendix~\ref{sec:baselines}.

\vspace*{-0.12in}
\paragraph{Results.}
Tables~\ref{tab:main1} and \ref{tab:main2} showcase comparative performance metrics for a diverse selection of datasets. For robustness, each experiment is repeated three times, varying the random split between training and test samples. The mean and standard deviation of the resulting AUC scores are then reported. \textsf{FeatLLM} demonstrates superior results, consistently finishing either first or second relative to the baseline methods under consideration. Remarkably, our approach delivers performance on par with traditional tabular models using 32 shots, despite employing only 4 shots (see Fig.~\ref{fig:compares}). Although methods such as STUNT and TabPFN yield notable performance gains on select datasets, their aggregate improvement over standard machine learning methods remains limited. Furthermore, models leveraging LLMs—including in-context learning, TABLET, and TabLLM—were unable to conduct inference on tabular tasks characterized by a high feature count. In contrast, our methodology, which integrates bagging and ensembling strategies, remains robust and effective across datasets with extensive feature spaces. It should also be acknowledged that extending our bagging and ensembling technique to other LLM-based frameworks is conceptually possible; however, this would essentially double per-sample inference operations, considerably escalating computational cost and thus rendering such an approach infeasible in practice.

\subsection{Analysis \& Discussion}
\label{sec:analysis}
\paragraph{Ablation study.} To assess the contribution of principal components, ablation experiments are conducted on the following elements: (1) \textsf{-Tuning}: exclusion of the weight-tuning phase in the linear model, relying instead on the summed binary feature for logit computation; (2) \textsf{-Ensemble}: elimination of the ensemble step; (3) \textsf{-Description}: removal of the feature description; and (4) \textsf{-Reasoning}: omission of the Step-1 instruction for rule generation.

Results of these ablations, averaged as AUC differences across all datasets, are presented in Table~\ref{tab:ablation}. Altering any of these components consistently results in diminished performance. Notably, the advantage afforded by weight tuning becomes more pronounced as the number of shots increases, suggesting that with larger datasets, precise estimation of rule weights is both feasible and impactful. Conversely, incorporating feature descriptions and detailed reasoning becomes most beneficial with limited shots, highlighting the importance of harnessing the inherent prior knowledge within LLMs for improving performance in few-shot contexts. Finally, across all experimental configurations, the ensemble strategy reliably bolsters predictive performance.

\vspace*{-0.12in}
\paragraph{Training \& inference time comparison.} We assess and contrast both training and inference durations across several methods, conducting our evaluation on the Adult dataset using a 16-shot training subset. Experiments default to a single A100 GPU, though TabLLM utilizes four A100 GPUs to enable model parallelism. Table~\ref{tab:cost} displays the measured runtimes. \textsf{FeatLLM} is notable for its relatively efficient inference phase, with timing that is on par with standard approaches. Conversely, LLM-based techniques—including In-context learning, TABLET, and TabLLM—demonstrate substantially slower inference speeds. Regarding training overhead, \textsf{FeatLLM} aligns closely with other LLM-oriented baselines: it only issues 30 API calls, resulting in minimal expense and potential for execution in parallel, thereby limiting any resource bottleneck.

\vspace*{-0.12in}
\paragraph{Quality of features generated by \textsf{FeatLLM}.}
To evaluate the utility of the rule-driven features produced by \textsf{FeatLLM} in addressing real-world tasks, we carry out a straightforward analysis. For each dataset, we calculate the average AUROC between newly derived features and target labels, subsequently determining what proportion of the generated features achieve an AUROC exceeding 0.5. These findings are organized in Table~\ref{tab:quality}. Results substantiate that most generated rules have substantial relevance to the target variable and yield significant predictive value for the underlying task (see column ``Before tuning"). In addition, the linear model fine-tuning process intrinsically filters out less relevant rules (i.e., those whose learned coefficients fall beneath a defined threshold), as seen in the column ``After tuning." The threshold is specified at 0.1, in accordance with the broader distribution of weight magnitudes.

\vspace*{-0.12in}
\paragraph{Effect of adding spurious correlations.} To evaluate the ability of different methods to eliminate irrelevant spurious correlations in scenarios with limited data, we performed a focused analysis. We experimented with the Heart dataset by artificially augmenting it with randomly selected columns from the Adult dataset that do not bear relevance to the Heart dataset’s objective. Model performance was then assessed as we incrementally increased the number of extraneous columns. To ensure an equitable comparison, all approaches were provided with an identical set of 8 labeled samples, and we omitted any methods (e.g., SCARF and STUNT) that require additional unlabeled data.

As illustrated in Figure~\ref{fig:spurious}, these added spurious correlations have varying impacts on model accuracy. \textsf{FeatLLM} experienced the smallest drop in accuracy compared to the baseline models. This resilience can be attributed to our method’s reliance on prior knowledge when extracting rules: it explicitly links features to the task, thereby disfavoring rules involving irrelevant attributes and conferring robustness to the model. Consistent with this, only a minor proportion (1–2\%) of rules derived from the noisy columns. Nevertheless, it is noteworthy that \textsf{FeatLLM} may still absorb spurious correlations and potentially conflate causality and correlation if augmented columns originate from datasets with attributes closely related to the Heart dataset (refer to Appendix~\ref{sec:spurious-appendix} for further analysis).

\vspace*{-0.12in}
\paragraph{Hyper-parameter analysis}
\textsf{FeatLLM} involves two primary hyper-parameters: the ensemble count and the number of rules retrieved per class during each LLM call. We systematically examined how these hyper-parameters affect model outcomes. Empirical results indicate that increasing the ensemble size consistently enhances performance up to approximately 20 ensembles, beyond which additional gains are marginal (refer to Fig.~\ref{fig:appendix-hyperparameter1} in the Appendix). With respect to the number of rules, although differences were not statistically significant, selecting 10 rules struck an optimal balance between prompt length and predictive accuracy (see Fig.~\ref{fig:hyperparameter2}). We attribute these findings to the fact that incorporating more rules increases input dimensionality, thereby introducing more information but also elevating the risk of overfitting, especially in low-shot settings.

\section{Conclusion}
We introduce \textsf{FeatLLM}, a method that advances large language model-based few-shot learning for tabular data. Unlike approaches that rely exclusively on LLMs for direct sample-wise predictions, \textsf{FeatLLM} employs LLMs as feature engineering tools, generating prediction criteria through rule induction. By leveraging rule creation in combination with ensemble techniques such as bagging, \textsf{FeatLLM} delivers substantial predictive performance with minimal inference costs. Notably, it relies solely on API access (no model training required) and can handle scenarios constrained by limited data features. Experimental results demonstrate that \textsf{FeatLLM} sets a new benchmark for data-efficient application of LLMs on real-world tabular datasets.

\textsf{FeatLLM} is currently optimized for situations with extremely limited training examples. In future research, we plan to extend its ability to construct novel features for datasets with more abundant samples, investigate the utility of alternative feature types beyond simple rules, and enhance interpretability to broaden applicability across diverse domains.

\section{Broader Impact}
The presented \textsf{FeatLLM} framework leverages the inherent knowledge within LLMs to achieve superior performance in tabular learning compared to standard supervised learning paradigms. By capitalizing on such prior knowledge, our approach makes key strides towards addressing overfitting challenges endemic to small sample training scenarios. \textsf{FeatLLM} offers significant advantages for practitioners, especially in domains like finance and healthcare, where data annotation is costly or otherwise restricted, and LLM pretraining on domain-relevant corpora can offer valuable supplemental insights. To foster broader adoption and safe deployment, it will be essential for future work to rigorously evaluate the potential influence of societal biases and misinformation embedded in the LLMs' prior knowledge, given their direct effect—and at times, dominance—over the model's outputs, regardless of the observed labeled data.

\end{document}